\documentclass[11pt,a4paper]{article}

\usepackage[utf8]{inputenc}
\usepackage[T1]{fontenc}
\usepackage[french]{babel}
\usepackage{amsmath}
\usepackage{amssymb}
\usepackage{listings}
\usepackage{xcolor}
\usepackage{geometry}
\usepackage{tikz-timing}
\geometry{hmargin=2.5cm, vmargin=2.5cm}

% Configuration de l'affichage du code VHDL
\lstset{
    language=VHDL,
    basicstyle=\ttfamily\small,
    keywordstyle=\color{blue}\bfseries,
    commentstyle=\color{green!50!black}\itshape,
    stringstyle=\color{red},
    numbers=left,
    numberstyle=\tiny\color{gray},
    stepnumber=1,
    frame=single,
    breaklines=true,
    backgroundcolor=\color{gray!5}
}

\title{\textbf{Correction d'Examen : Architecture des Ordinateurs}}
\author{Grenoble INP - Esisar \cite{1170} \\ Cours CE317 \cite{1173}}
\date{Session 1 - Année 2015-2016 \cite{1175} (Durée : 180mn \cite{1172})}

\begin{document}

\maketitle

\section*{Partie I : Partie cours [34/90] \cite{1183}}

\subsection*{Question 1.1 \cite{1184}}
\textbf{Énoncé :} Codez en virgule fixe avec le minimum de bits 511,75 [6] \cite{1184}.

\textbf{Correction et Explications :}
Pour coder un nombre en virgule fixe, on convertit séparément la partie entière et la partie décimale en binaire.
\begin{itemize}
    \item \textbf{Partie entière (511)} : On sait que $2^9 = 512$, donc $511 = 2^9 - 1$. En binaire, 511 s'écrit avec 9 bits à '1' : $111111111_2$.
    \item \textbf{Partie décimale (0,75)} : $0,75 = 0,5 + 0,25 = 2^{-1} + 2^{-2}$. En binaire, cela s'écrit $0,11_2$.
\end{itemize}
\textbf{Réponse :}
\begin{itemize}
    \item En \textbf{non signé}, le minimum est de \textbf{11 bits} : $111111111,11_2$.
    \item En \textbf{signé} (complément à 2), il faut ajouter un bit de signe (0 pour positif), soit un minimum de \textbf{12 bits} : $0111111111,11_2$.
\end{itemize}

\subsection*{Question 1.2 \cite{1185, 1186}}
\textbf{Énoncé :} Simplifiez en utilisant un tableau de Karnaugh la fonction $F$ [12] \cite{1185}.

\textit{Note : Le texte d'origine a perdu les barres de négation (OCR : F=ABC+ABC+ABC+ABC+ABC \cite{1186}). Nous traiterons l'exemple suivant pour illustrer la méthode : $F = \overline{A}\overline{B}C + \overline{A}BC + A\overline{B}C + AB\overline{C} + ABC$}

\textbf{Correction et Explications :}
On place des '1' dans le tableau de Karnaugh pour chaque minterme :
\begin{center}
\begin{tabular}{|c|c|c|c|c|}
\hline
$A \backslash BC$ & 00 & 01 & 11 & 10 \\
\hline
0 & 0 & \textbf{1} ($\overline{A}\overline{B}C$) & \textbf{1} ($\overline{A}BC$) & 0 \\
\hline
1 & 0 & \textbf{1} ($A\overline{B}C$) & \textbf{1} ($ABC$) & \textbf{1} ($AB\overline{C}$) \\
\hline
\end{tabular}
\end{center}
Regroupements possibles (les plus grands possibles en puissances de 2) :
\begin{itemize}
    \item Un groupe de 4 (colonne 01 et 11) : donne $C$.
    \item Un groupe de 2 (ligne du bas, 11 et 10) : donne $AB$.
\end{itemize}
\textbf{Résultat simplifié :} $F = C + AB$.

\subsection*{Question 1.3 \cite{1187, 1188, 1189}}
\textbf{Énoncé :} En VHDL, quel type normalisé doit être utilisé pour réaliser des opérations arithmétiques sur des vecteurs binaires non signés [2] \cite{1187} ? Quels bibliothèques et paquetages déclarer [2] \cite{1188} ? Montrez comment $-5 - 7$ est obtenu en complément à 2 \cite{1189}.

\textbf{Correction et Explications :}
\begin{itemize}
    \item \textbf{Type normalisé :} Il faut utiliser le type \texttt{unsigned} \cite{1187}.
    \item \textbf{Bibliothèque et paquetage :} \texttt{library IEEE; use IEEE.numeric\_std.all;} \cite{1188} (ne jamais utiliser l'obsolète \texttt{std\_logic\_unsigned}).
    \item \textbf{Opération $-5 - 7$ en complément à 2 (sur 5 bits) :} \cite{1189}
    Pour coder $-12$ (le résultat), il faut au minimum 5 bits (plage de $-16$ à $+15$).
    \begin{enumerate}
        \item $+5$ en binaire = $00101_2$. Inversion = $11010_2$. On ajoute 1 $\rightarrow$ \textbf{$-5 = 11011_2$}.
        \item $+7$ en binaire = $00111_2$. Inversion = $11000_2$. On ajoute 1 $\rightarrow$ \textbf{$-7 = 11001_2$}.
        \item Addition : $11011_2 + 11001_2 = (1) 10100_2$. On ignore la retenue sortante.
        \item Résultat : \textbf{$10100_2$}. Pour vérifier sa valeur : on inverse ($01011_2$) et on ajoute 1 ($01100_2 = 12_{10}$). Le résultat est bien $-12$.
    \end{enumerate}
\end{itemize}

\subsection*{Question 1.4 \cite{1190, 1191}}
\textbf{Énoncé :} En VHDL est-ce que l'ordre des processus a une importance \cite{1190} ? Si non, dans quel ordre sont-ils exécutés \cite{1191} ?

\textbf{Correction et Explications :}
\textbf{Non}, l'ordre des processus dans l'architecture n'a aucune importance \cite{1190}. 
En VHDL, les processus sont des instructions \textbf{concurrentes}. Ils s'exécutent de manière parallèle \cite{1191}. Leur exécution est déclenchée (réveillée) uniquement par les événements qui surviennent sur les signaux présents dans leur \textbf{liste de sensibilité}.

\newpage

\section*{Partie II : Conception d'un registre à décalage SIPO [56/90] \cite{1199, 1200}}

\subsection*{Question 2.1 \cite{1201}}
\textbf{Énoncé :} Dessinez le schéma d'un registre SIPO de $n$ bits en utilisant $n$ bascules D \cite{1201}. Entrées : \texttt{din}, \texttt{clk}, \texttt{enable}, \texttt{reset} \cite{1203, 1204, 1205, 1207}. Sortie : \texttt{dout} (n bits) \cite{1206}.

\textit{(Voir le code SVG fourni séparément pour générer ou visualiser ce schéma).} Explication : Les bascules sont en cascade. L'entrée $D$ de la première bascule est liée à \texttt{din}. L'entrée $D$ des bascules suivantes est liée à la sortie $Q$ de la bascule précédente. \texttt{clk}, \texttt{reset} et \texttt{enable} sont communs à toutes les bascules.

\subsection*{Question 2.2 \cite{1214}}
\textbf{Énoncé :} Ecrire en VHDL RTL l'entité générique (\texttt{shift}) et l'architecture (\texttt{rtl}) \cite{1214}.

\textbf{Correction :}
\begin{lstlisting}
library IEEE;
use IEEE.std_logic_1164.all;

entity shift is
    generic (n : integer := 8);
    port (
        din    : in  std_logic;
        clk    : in  std_logic;
        reset  : in  std_logic;
        enable : in  std_logic;
        dout   : out std_logic_vector(n-1 downto 0)
    );
end entity shift;

architecture rtl of shift is
    signal shift_reg : std_logic_vector(n-1 downto 0);
begin
    process(clk)
    begin
        if rising_edge(clk) then
            if reset = '1' then
                shift_reg <= (others => '0');
            elsif enable = '1' then
                -- Decalage vers la gauche : le MSB recoit din
                shift_reg <= shift_reg(n-2 downto 0) & din;
            end if;
        end if;
    end process;
    dout <= shift_reg;
end architecture rtl;
\end{lstlisting}

\subsection*{Question 2.3 \cite{1215}}
\textbf{Énoncé :} Ecrire en VHDL structurel l'architecture (\texttt{struct}) de la même entité \cite{1215}. Composant disponible : \texttt{bascule} \cite{1210-1213}.

\textbf{Correction :}
\begin{lstlisting}
architecture struct of shift is
    component bascule
        port(d, clk, reset, enable: in std_logic;
             q: out std_logic);
    end component;
    
    -- Signal intermediaire pour relier les bascules en cascade
    signal q_int : std_logic_vector(n downto 0);
begin
    q_int(0) <= din; -- L'entree du registre attaque la 1ere bascule

    gen_shift: for i in 0 to n-1 generate
        bascule_inst : bascule port map (
            d      => q_int(i),
            clk    => clk,
            reset  => reset,
            enable => enable,
            q      => q_int(i+1)
        );
    end generate gen_shift;

    -- Affectation des sorties paralleles
    dout <= q_int(n downto 1);
end architecture struct;
\end{lstlisting}

\subsection*{Question 2.4 \cite{1216}}
\textbf{Énoncé :} Ecrire un testbench simulant l'architecture structurelle pour $n=8$ \cite{1216, 1217}.

\textbf{Correction :}
\begin{lstlisting}
library IEEE;
use IEEE.std_logic_1164.all;

entity tb_shift is
end entity tb_shift;

architecture sim of tb_shift is
    signal din, clk, reset, enable : std_logic := '0';
    signal dout : std_logic_vector(7 downto 0);
    constant period : time := 10 ns;
begin
    UUT: entity work.shift(struct)
        generic map (n => 8)
        port map (din => din, clk => clk, reset => reset, enable => enable, dout => dout);

    clk_proc: process begin
        clk <= '0'; wait for period/2;
        clk <= '1'; wait for period/2;
    end process;

    stim_proc: process begin
        reset <= '1'; wait for 20 ns;
        reset <= '0'; enable <= '1'; din <= '1'; wait for period;
        din <= '0'; wait for 2 * period;
        enable <= '0'; wait for 2 * period; -- Retenue d'etat
        enable <= '1'; din <= '1'; wait for 2 * period;
        reset <= '1'; wait for period;
        wait;
    end process;
end architecture sim;
\end{lstlisting}

\subsection*{Question 2.5 \cite{1218}}
\textbf{Énoncé :} Dessinez les chronogrammes correspondant au testbench \cite{1218, 1219}.

\textbf{Correction (via TikZ-Timing) :}
\begin{center}
\begin{tikztimingtable}[timing/slope=0]
  \texttt{clk}    & 14{C} \\
  \texttt{reset}  & 2{1} 10{0} 2{1} \\
  \texttt{enable} & 2{0} 6{1} 2{0} 4{1} \\
  \texttt{din}    & 2{0} 2{1} 6{0} 4{1} \\
  \texttt{dout}   & 2D{00} 2D{00} 2D{01} 2D{02} 2D{04} 2D{04} 2D{09} 2D{00} \\
\extracode
  \tablegrid
\end{tikztimingtable}
\end{center}

\end{document}